\subsection{Cogebra Structure on the Restricted Dual}

For any \(a\) in \(\mathbf{A}\) , we denote by \(\eta(a)\) the linear form \(f \mapsto f(a)\) on \(\Theta(\mathbf{A})\) . We thus define a K-linear mapping \(\eta\) from \(\mathbf{A}\) to the dual \(\Theta(\mathbf{A})^*\) of the vector space \(\Theta(\mathbf{A})\) . Set \(\varepsilon = \eta(1)\) .

PROPOSITION 8. --- On the vector space \(\Theta(A)\) , there exists a unique cogebra structure (III, §11, No.~1, p.~574) such that the mapping \(\eta: A \to \Theta(A)^*\) is a homomorphism from \(A\) to the dual algebra (III, §11, No.~2, p.~579) of \(\Theta(A)\) . The cogebra \(\Theta(A)\) is coassociative and admits \(\varepsilon\) as a counit.

For every integer \(n \geqslant 1\) , consider the K-linear mapping \(j_{n}\) from \(\Theta(\mathrm{A})^{\otimes n}\) to the dual of \(A^{\otimes n}\) characterized by the formula

\[
\langle j _ {n} \left(f _ {1} \otimes \dots \otimes f _ {n}\right), a _ {1} \otimes \dots \otimes a _ {n} \rangle = \prod_ {i = 1} ^ {n} \left\langle f _ {i}, a _ {i} \right\rangle\tag{13}
\]

for \((a_{i})\) in \(\mathrm{A}^n\) and \((f_i)\) in \(\Theta (\mathrm{A})^n\) . By Proposition 16, (ii) of II, §7, No.~7, p.~308, the mapping \(j_{n}\) is injective. Denote by \(m_{\mathrm{K}}:\mathrm{K}\otimes \mathrm{K}\to \mathrm{K}\) and \(m_{\mathrm{A}}:\mathrm{A}\otimes \mathrm{A}\rightarrow \mathrm{A}\) the mappings deduced from the multiplication in K and A, respectively. For \(f,g\in \Theta (\mathrm{A})\) and \(a,b\in \mathrm{A}\) , we have

\[
\langle j _ {2} (f \otimes g), a \otimes b \rangle = m _ {\mathrm{K}} \circ (\eta (a) \otimes \eta (b)) (f \otimes g).
\]

We therefore have

\[
\langle j _ {2} (t), a \otimes b \rangle = m _ {\mathrm{K}} \circ (\eta (a) \otimes \eta (b)) (t)\tag{14}
\]

for every \(t\in \Theta (\mathrm{A})\otimes \Theta (\mathrm{A})\).

Lemma 2. --- Let \(c: \Theta(A) \to \Theta(A) \otimes \Theta(A)\) be a K-linear mapping. Then \(\eta\) is a homomorphism from A to the dual algebra of the cogebra \((\Theta(A), c)\) if and only if the following diagram commutes:

\[
\begin{array}{c} \Theta (\mathrm{A}) \xrightarrow {c} \Theta (\mathrm{A}) \otimes \Theta (\mathrm{A}) \\ \Biggl \downarrow_ {j _ {1}} \qquad \qquad \qquad \qquad \Biggl \downarrow_ {j _ {2}} \\ \mathrm{A} ^ {*} \xrightarrow {^ t m _ {\mathrm{A}}} (\mathrm{A} \otimes \mathrm{A}) ^ {*}. \end{array}\tag{15}
\]

Indeed, \(\eta\) is a homomorphism from A to the dual algebra of the cogebra \((\Theta (\mathrm{A}),c)\) if and only if we have \(\eta (ab) = m_{\mathrm{K}}\circ (\eta (a)\otimes \eta (b))\circ c\) for all \(a,b\in \mathrm{A}\) , that is

\[
\eta (a b) (f) = m _ {\mathrm{K}} \circ (\eta (a) \otimes \eta (b)) (c (f))
\]

for \(a, b \in \mathrm{A}\) and \(f \in \Theta(\mathrm{A})\) . Now, we have

\[
\eta (a b) (f) = f (a b) = \left\langle {} ^ {t} m _ {\mathrm{A}} \left(j _ {1} (f)\right), a \otimes b \right\rangle
\]

and, by (14),

\[
m _ {K} \circ (\eta (a) \otimes \eta (b)) (c (f)) = \langle j _ {2} (c (f)), a \otimes b \rangle
\]

for all \(a, b \in \mathrm{A}\) and every \(f \in \Theta(\mathrm{A})\) . The lemma follows.

Since \(j_{2}\) is injective, there exists at most one linear mapping c that makes the diagram above commute. To prove its existence, we must prove that the image of \(^{t}m \circ j_{1}\) is contained in that of \(j_{2}\) . In other words, we must prove that there exist, for every element f of \(\Theta(\mathrm{A})\) , a natural number n and elements \(f_{1}^{\prime}, \ldots, f_{n}^{\prime}, f_{1}^{\prime\prime}, \ldots, f_{n}^{\prime\prime}\) of \(\Theta(\mathrm{A})\) satisfying the relations

\[
f (a b) = \sum_ {i = 1} ^ {n} f _ {i} ^ {\prime} (a) f _ {i} ^ {\prime \prime} (b)\tag{16}
\]

for \(a, b \in \mathbf{A}\) . We will then have

\[
c (f) = \sum_ {i = 1} ^ {n} f _ {i} ^ {\prime} \otimes f _ {i} ^ {\prime \prime}.\tag{17}
\]

By the corollary of VIII, p.~379, there exists a left A-module E of finite dimension over K with f as a coefficient. Let \((e_{1},\ldots,e_{n})\) be a basis of E, \((e_{1}^{*},\ldots,e_{n}^{*})\) the dual basis, x an element of E, and \(x^{*}\) an element of \(E^{*}\) such that \(f=c_{\mathrm{E}}(x,x^{*})\) . Set \(f_{i}^{\prime}=c_{\mathrm{E}}(e_{i},x^{*})\) and \(f_{i}^{\prime\prime}=c_{\mathrm{E}}(x,e_{i}^{*})\) for \(i\in[1,n]\) ; for

\(a, b\) in A, we have

\[
\begin{array}{r l} & {f (a b) = \langle x ^ {*}, a b x \rangle = \langle x ^ {*} a, b x \rangle = \left\langle \sum_ {i} \langle x ^ {*} a, e _ {i} \rangle e _ {i} ^ {*}, b x \right\rangle} \\ & {\qquad = \sum_ {i} \langle x ^ {*} a, e _ {i} \rangle \langle e _ {i} ^ {*}, b x \rangle = \sum_ {i} \langle x ^ {*}, a e _ {i} \rangle \langle e _ {i} ^ {*}, b x \rangle = \sum_ {i} f _ {i} ^ {\prime} (a) f _ {i} ^ {\prime \prime} (b),} \end{array}
\]

and therefore (16).

Let us prove the coassociativity of \(c\) . For this, consider the K-linear mappings

\[
c ^ {\prime} = \left(c \otimes 1 _ {\Theta (\mathrm{A})}\right) \circ c \quad \text { and } \quad c ^ {\prime \prime} = \left(1 _ {\Theta (\mathrm{A})} \otimes c\right) \circ c
\]

from \(\Theta (\mathrm{A})\) to \(\Theta (\mathrm{A})^{\otimes 3}\) . We have the relations

\[
\begin{array}{r l} \langle j _ {3} (f \otimes c (g)), a \otimes b \otimes c \rangle & = \langle f, a \rangle \langle j _ {2} \circ c (g), b \otimes c \rangle \\ & = \langle f, a \rangle \langle g, b c \rangle \\ & = \langle j _ {2} (f \otimes g), a \otimes b c \rangle \\ & = \langle^ {t} (\mathrm{Id} _ {\mathrm{A}} \otimes m _ {\mathrm{A}}) \circ j _ {2} (f \otimes g), a \otimes b \otimes c \rangle \end{array}
\]

for \(f, g \in \Theta(\mathbf{A})\) and \(a, b, c \in \mathbf{A}\) . From this, we deduce that the following diagram commutes:

\[
\begin{array}{c} \Theta (\mathrm{A}) \otimes \Theta (\mathrm{A}) \xrightarrow {\operatorname{Id} _ {\Theta (\mathrm{A})} \otimes c} \Theta (\mathrm{A}) \otimes \Theta (\mathrm{A}) \otimes \Theta (\mathrm{A}) \\ \Biggl \downarrow_ {j _ {2}} \\ (\mathrm{A} \otimes \mathrm{A}) ^ {*} \xrightarrow {^ t (\operatorname{Id} _ {\mathrm{A}} \otimes m _ {\mathrm{A}})} (\mathrm{A} \otimes \mathrm{A} \otimes \mathrm{A}) ^ {*}. \end{array}\tag{18}
\]

Because of the commutativity of this diagram and that of (15), for \(f \in \Theta(\mathrm{A})\) and \(a, a', a'' \in A\) , we have

\[
\langle j _ {3} \circ c ^ {\prime} (f), a \otimes a ^ {\prime} \otimes a ^ {\prime \prime} \rangle = \langle f, (a a ^ {\prime}) a ^ {\prime \prime} \rangle ;
\]

we can show the relation

\[
\langle j _ {3} \circ c ^ {\prime \prime} (f), a \otimes a ^ {\prime} \otimes a ^ {\prime \prime} \rangle = \langle f, a (a ^ {\prime} a ^ {\prime \prime}) \rangle
\]

likewise. Since multiplication in A is associative, we have \(j_{3} \circ c' = j_{3} \circ c''\) , and therefore \(c' = c''\) because \(j_{3}\) is injective.

Finally, formulas (16) and (17) imply that \(\Theta(\mathrm{A})\) admits \(\varepsilon\) as a counit.

Remark 1. --- Let \((\mathrm{V},\pi)\) be a finite-dimensional linear representation of the algebra A. Let us introduce a basis \((e_{1},\ldots,e_{n})\) of V and the dual basis \((e_{1}^{*},\ldots,e_{n}^{*})\) of \(V^{*}\) . By the proof of Lemma 2, we have the relation

\[
c (c _ {\pi} (x, x ^ {*})) = \sum_ {k = 1} ^ {n} c _ {\pi} (e _ {k}, x ^ {*}) \otimes c _ {\pi} (x, e _ {k} ^ {*})\tag{19}
\]

for \(x \in V\) and \(x^{*} \in V^{*}\) . For \(1 \leqslant i, j \leqslant n\) , set \(\pi_{ij} = c_{\pi}(e_j, e_i^*)\) . For every \(a \in A\) , the matrix of \(\pi(a)\) with respect to the basis \((e_1, \ldots, e_n)\) of V is equal to \((\pi_{ij}(a))\) . For \(1 \leqslant i \leqslant n\) and \(1 \leqslant j \leqslant n\) , we then have

\[
c (\pi_ {i j}) = \sum_ {k = 1} ^ {n} \pi_ {i k} \otimes \pi_ {k j}.\tag{20}
\]

DEFINITION 3. --- Let C be a cogebra over the field K and c its coproduct. A subcogebra of C is any linear subspace \(C_{1}\) of C such that \(c(\mathrm{C}_{1})\) is contained in the canonical image of \(C_{1} \otimes_{K} C_{1}\) in \(C \otimes_{K} C\) .

Let \(j\) be the canonical injection of \(C_1\) into \(C\) . By this definition, there exists a unique linear mapping \(c_1: C_1 \to C_1 \otimes_K C_1\) such that we have

\[
c \circ j = (j \otimes j) \circ c _ {1};\tag{21}
\]

this relation means that \(j\) is a morphism of cogebras from \((\mathrm{C}_1, c_1)\) to \((\mathrm{C}, c)\) (III, §11, No.~1, p.~574).

If C is coassociative, then \(C_{1}\) is coassociative. If C is cocommutative, then so is \(C_{1}\) . If C has a counit \(\varepsilon\) , then the restriction of \(\varepsilon\) to \(C_{1}\) is a counit of \(C_{1}\) .

PROPOSITION 9. --- Let \(\Theta\) be a linear subspace of \(\Theta(A)\) . The following properties are equivalent:

\begin{enumerate}
\def\labelenumi{(\roman{enumi})}
\item
  \(\Theta\) is an (A, A)-sub-bimodule of \(\Theta(A)\) .
\item
  \(\Theta\) is a subcogebra of \(\Theta(A)\) .
\item
  There exists a hereditary set \(\mathcal{C}\) of classes of A-modules of finite dimension over \(K\) such that \(\Theta = \Theta_{\mathcal{C}}(A)\) .
\end{enumerate}

When these properties hold, the set C mentioned in (iii) is uniquely determined.

The last assertion follows from the corollary of VIII, p.~388: the set C consists of the classes of A-modules E of finite dimension over K such that \(\Theta_{\mathrm{E}}(\mathrm{A})\) is contained in \(\Theta\) .

\begin{enumerate}
\def\labelenumi{(\roman{enumi})}
\setcounter{enumi}{2}
\item
  \(\Rightarrow\) (ii): Let C be a hereditary set of classes of A-modules of finite dimension over K. Then \(\Theta_{\mathcal{C}}(\mathrm{A})\) is the union of the directed family \((\Theta_{\mathrm{E}}(\mathrm{A}))_{\mathrm{E}\in\mathcal{C}}\) . Since \(\Theta_{\mathrm{E}}(\mathrm{A})\) is a subcogebra of \(\Theta(\mathrm{A})\) for every \(E\in C\) (VIII, p.~390, formula (19)), the same holds for \(\Theta_{\mathcal{C}}(\mathrm{A})\) .
\item
  \(\Rightarrow\) (i): Let \(f\in \Theta (\mathrm{A})\) . Let \(f_{1}^{\prime},\ldots ,f_{n}^{\prime},f_{1}^{\prime \prime},\ldots ,f_{n}^{\prime \prime}\) be elements of \(\Theta (\mathrm{A})\) satisfying \(c(f) = \sum f_i^\prime \otimes f_i''\) . For \(a,b\) in A, we have \(f(ab) = \sum f_i'(a)f_i''(b)\) , and therefore
\end{enumerate}

\[
b f = \sum_ {i = 1} ^ {n} f _ {i} ^ {\prime \prime} (b) f _ {i} ^ {\prime}, \quad f a = \sum_ {i = 1} ^ {n} f _ {i} ^ {\prime} (a) f _ {i} ^ {\prime \prime}
\]

(VIII, p.~375, formulas (3) and (4)). Consequently, a subcogebra of \(\Theta(A)\) is an \((A, A)\) -sub-bimodule.

\begin{enumerate}
\def\labelenumi{(\roman{enumi})}
\tightlist
\item
  \(\Rightarrow\) (iii): Suppose that \(\Theta\) is an (A, A)-sub-bimodule of \(\Theta(\mathrm{A})\) ; let C be the set of classes of A-modules E of finite dimension over K such that \(\Theta_{\mathrm{E}}(\mathrm{A})\) is contained in \(\Theta\) . The set C is hereditary (VIII, p.~377, Remark 2), and we have \(\Theta_{\mathcal{C}}(\mathrm{A}) \subset \Theta\) by construction. Let \(f \in \Theta\) and E = Af. Then E is finite-dimensional over K. Consequently, every linear form on E is of the form \(u_{a}: g \mapsto g(a)\) with a in A (II, §7, No.~5, p.~302, Corollary 2). Now, for \(a, b, x \in A\) , we have
\end{enumerate}

\[
c _ {\mathrm{E}} (a f, u _ {b}) (x) = \langle u _ {b}, x a f \rangle = f (b x a) = a f b (x),
\]

so that \(c_{\mathrm{E}}(af, u_{b}) = afb\) . We therefore have \(\Theta_{\mathrm{E}}(\mathrm{A}) \subset \Theta\) . Consequently, the A-module E is of type C, and f is one of its coefficients. We therefore have \(\Theta \subset \Theta_{\mathcal{C}}(\mathrm{A})\) and, finally, \(\Theta = \Theta_{\mathcal{C}}(\mathrm{A})\) .

Remark 2. --- Let \(\Theta\) be a subcogebra of \(\Theta(A)\) . We endow \(K\) with the discrete topology and the algebra \(A\) with the coarsest topology for which the mappings \(f: A \to K\) for \(f\) running through \(\Theta\) are continuous (Gen.~Top., I, §2, No.~2, p.~175). This topology endows \(A\) with the structure of a topological \(K\) -module. A two-sided ideal \(\mathfrak{a}\) of \(A\) is open if and only if it has finite codimension and its orthogonal \(\mathfrak{a}'\) is contained in \(\Theta\) . By Proposition 3 of VIII, p.~379, the open two-sided ideals of \(A\) form a fundamental system of neighborhoods of \(0\) in \(A\) . The topology on \(A\) is therefore compatible with its ring structure.

Let C be the hereditary set of classes of A-modules of finite dimension over K such that \(\Theta = \Theta_{C}\) (VIII, p.~391, Proposition 9). Let E be a left A-module of finite dimension over K. We endow it with the discrete topology. The following properties are equivalent:

\begin{enumerate}
\def\labelenumi{(\roman{enumi})}
\item
  The A-module E is of type C.
\item
  The annihilator of the A-module E is open in A.
\item
  The mapping \((a, x) \mapsto ax\) from A × E to E is continuous.
\end{enumerate}

The last property means that \(\mathrm{E}\) is a topological A-module.

Let \(\Theta^{*}\) be the dual algebra of the cogebra \(\Theta\) . We endow \(\Theta^{*}\) with the coarsest topology for which the mappings \(\varphi \mapsto \varphi(u)\) from \(\Theta^{*}\) to K, for \(u\) running through \(\Theta\) , are continuous. The topology on the algebra \(\Theta^{*}\) is compatible with the additive group structure on \(\Theta^{*}\) . The orthogonals in \(\Theta^{*}\) of the sets of the form \(\Theta_{\mathrm{E}}(\mathrm{A})\) , where E is an A-module of type \(\mathcal{C}\) , form a fundamental system of neighborhoods of 0. Now, such a set is a subcogebra of \(\Theta\) , so its orthogonal is an ideal of \(\Theta^{*}\) . The topology on \(\Theta^{*}\) is therefore compatible with its ring structure. The canonical algebra homomorphism \(\eta: A \to \Theta^{*}\) (that sends \(a \in A\) to the linear form \(f \mapsto f(a)\) on \(\Theta\) ) defines an isomorphism from the completed Hausdorff space \(\widehat{A}\) of A (Gen.~Top., II, §3, No.~7, p.~191) to \(\Theta^{*}\) .

